\section{Future Work}

The first work to do next is to write a Pit compiler in Pit.  This will both demonstrate that Pit can be self-hosting
and improve on the current compiler, which has an overly simple design with no optimizations.  It was written for a
single purpose, which is to bootstrap the compiler written in Pit.

Then the compiler needs some optimization.  Most of the optimization techniques used by C compilers should be useful in
Pit.  There are additional Pit-specific optimizations.  For example, a primitive variable can often be substituted for
an auto variable as described above.  When work on optimization is complete, it is reasonable to expect Pit programs
that use only primitive variables to match the speed of equivalent C programs.  Pit programs should always be faster
than equivalent Perl programs because when a program written in Perl executes, Perl first compiles it into byte-code,
then executes it in a virtual machine.  Pit programs should execute significantly faster, even if they only use auto
variables, because the code is already compiled at run-time and does not require a virtual machine.  This places the
performance of Pit programs, at worst, somewhere close to C programs and better than Perl programs.

Next, there needs to be some study of the difference in speed between primitive and auto variables, specifically to
determine how high the performance penalty is for using auto variables.  This will provide insight into when it's appropriate
to use primitive or auto variables when weighing the trade-off of a program's performance versus its security,
portability, and difficulty to write.  Some rough study on this could be done today on toy examples by emulating
the logic for auto variables in C.  The examples would be lengthy and inconvenient, but with care they could serve the
purpose of giving a rough estimate of the difference in speed between the types of variables.

A later useful branch of work may be to create a compiler that translates Pit code into a byte-code that is suitable
for a virtual machine.  Like Java, this would allow programs to be portable between CPU architectures.  It should be
possible to do this in a way that allows many programs to be compiled natively or for the virtual machine without
modification.

Another possible branch of work may be adapting existing research on checking C code for security bugs.  While this is
obviated for code that exclusively uses auto variables, this research may still be useful to minimize security risks in
code that uses primitive variables.
